\documentclass[12pt]{article}
% Préambule commun à tous les documents extraits de « La Revue CAPES »
\usepackage[T1]{fontenc}
\usepackage[utf8]{inputenc}
\usepackage{lmodern}
\usepackage{amsmath,amssymb,amsthm,amsfonts,mathrsfs}
\usepackage{setspace}
\usepackage{listings}
\usepackage{pgfplots}
\pgfplotsset{compat=1.18}
\usepackage{longtable}
\usepackage{adjustbox}
\usepackage{lettrine}
\usepackage{tkz-tab}
\usepackage{changepage}
\usepackage{pdflscape}
\usepackage{graphicx}
\usepackage{tikz}
\usepackage{xcolor}
\usepackage[a4paper,top=2cm,bottom=2.5cm,left=2.5cm,right=2cm]{geometry}
\usepackage{fancyhdr}
\usepackage{draftwatermark}
\SetWatermarkText{La RevueCapes}
\SetWatermarkLightness{0.9}
\SetWatermarkScale{2}
\usepackage{hyperref}
\hypersetup{colorlinks=true,linkcolor=blue!50!black,urlcolor=blue!50!black}

\definecolor{revue}{RGB}{20,60,120}

\pagestyle{fancy}
\fancyhf{}
\renewcommand{\headrulewidth}{0.4pt}
\setlength{\headheight}{15pt}

% \entete{Matière}{Titre}{Type de document}
\newcommand{\entete}[3]{%
  \fancyhead[L]{\small\textsc{La Revue CAPES} -- #1}%
  \fancyhead[R]{\small #2}%
  \fancyfoot[C]{\thepage}%
  \thispagestyle{plain}%
  \begin{center}
    {\color{revue}\rule{\textwidth}{1.2pt}}\\[4pt]
    {\small\textsc{Concours directs CAP-CEG / CAPES -- Burkina Faso}}\\[6pt]
    {\LARGE\bfseries\color{revue} #2}\\[6pt]
    {\large\itshape #1 -- #3}\\[2pt]
    {\color{revue}\rule{\textwidth}{1.2pt}}
  \end{center}
  \vspace{0.5cm}%
}

% Encadré pour les remarques ajoutées dans les corrigés
\newcommand{\remarque}[1]{\par\medskip\noindent\fcolorbox{revue}{revue!6}{\parbox{\dimexpr\linewidth-2\fboxsep-2\fboxrule}{\textbf{\color{revue}Remarque.} #1}}\par\medskip}


\begin{document}
\entete{Culture générale}{Corrigé -- Session 2016}{Correction}
\begin{spacing}{1.5}

\noindent\textbf{Sujet.} Faisant écho à l'avènement de la société de l'information, Gaston Mialaret (\emph{Pédagogie générale}, 1991, p.~530) affirme : « \emph{avant on savait tout sur le village mais rien sur le monde ; aujourd'hui on sait tout sur le monde mais rien sur le village} ». Justifiez cette réflexion de Mialaret, présentez deux (02) conséquences possibles de cette situation et avancez deux (02) mesures que l'on pourrait mettre en œuvre pour essayer de trouver un équilibre.

\remarque{Ce sujet ne figurait pas parmi les corrigés du document d'origine ; la proposition ci-dessous a été rédigée pour compléter la série.}

\rubrique{1. Analyse du sujet}

\begin{itemize}
\item \textbf{Type de sujet :} la consigne impose trois tâches, donc un plan en trois parties : I. \textbf{justifier} la réflexion (montrer qu'elle est fondée) ; II. présenter \textbf{deux conséquences} ; III. proposer \textbf{deux mesures}. Il ne s'agit pas de la contester.
\item \textbf{L'auteur :} Gaston Mialaret (1918-2016), pédagogue français, spécialiste des sciences de l'éducation.
\item \textbf{Mots-clés :}
  \begin{itemize}
  \item « société de l'information » : société où l'information circule massivement et instantanément grâce aux médias (radio, télévision) puis à Internet, aux téléphones portables et aux réseaux sociaux ;
  \item « le village » : le milieu local, la communauté d'origine, ses traditions, son histoire, ses langues, ses réalités ;
  \item « le monde » : l'environnement planétaire, les événements et les cultures lointaines ;
  \item l'opposition « avant / aujourd'hui » et le chiasme « tout sur le village / rien sur le monde » $\to$ « tout sur le monde / rien sur le village » traduisent un renversement.
  \end{itemize}
\item \textbf{Problématique :} comment l'ouverture au monde permise par les médias s'est-elle accompagnée d'une méconnaissance du milieu local, quelles en sont les conséquences, et comment concilier enracinement et ouverture ?
\end{itemize}

\rubrique{2. Plan détaillé}

\begin{enumerate}
\item[I.] \textbf{Justification de la réflexion.} 1. Hier, une connaissance profonde du village mais un horizon limité. 2. Aujourd'hui, une ouverture au monde par les médias et le numérique. 3. Une méconnaissance croissante du milieu local (école, exode, médias tournés vers l'extérieur).
\item[II.] \textbf{Deux conséquences.} 1. L'acculturation et la perte des repères identitaires. 2. Le déracinement et le désintérêt pour le développement local.
\item[III.] \textbf{Deux mesures pour un équilibre.} 1. Une école qui enseigne aussi le milieu local et les langues nationales. 2. Valoriser la culture locale par les médias, le numérique et la transmission intergénérationnelle.
\end{enumerate}

\rubrique{3. Proposition de rédaction}

\partie{Introduction}

L'humanité vit depuis quelques décennies une révolution de la communication : la radio, la télévision, puis Internet, le téléphone portable et les réseaux sociaux mettent instantanément chacun en contact avec le reste de la planète. Faisant écho à l'avènement de cette société de l'information, le pédagogue Gaston Mialaret affirme dans \emph{Pédagogie générale} (1991) : « avant on savait tout sur le village mais rien sur le monde ; aujourd'hui on sait tout sur le monde mais rien sur le village ». Par cette formule construite comme un renversement, il souligne que l'ouverture sur le monde s'est accompagnée d'un éloignement à l'égard du milieu de vie immédiat. Comment justifier ce constat ? Quelles en sont les conséquences, et comment rétablir un équilibre entre la connaissance du monde et celle du village ? Nous justifierons d'abord la réflexion de Mialaret, puis nous présenterons deux conséquences de cette situation, avant de proposer deux mesures pour trouver un équilibre.

\partie{I. Justification de la réflexion de Mialaret}

\souspartie{1. Hier, une connaissance profonde du village mais un horizon limité}
Dans les sociétés traditionnelles, l'enfant était éduqué au sein de sa famille et de sa communauté. Il apprenait l'histoire de son lignage, les contes, les proverbes, les interdits, les techniques agricoles, la connaissance des plantes et des saisons, grâce aux veillées autour du feu et à l'enseignement des anciens. Il connaissait parfaitement son village. En revanche, faute de moyens de communication, il ignorait presque tout de ce qui se passait au-delà de sa région.

\souspartie{2. Aujourd'hui, une ouverture au monde}
Grâce aux médias et au numérique, un jeune de Ouagadougou ou de Fada N'Gourma peut suivre en direct un match de football européen, connaître les dernières nouvelles des États-Unis ou les tendances de la mode asiatique. L'école, elle aussi, enseigne la géographie et l'histoire du monde. Les jeunes connaissent ainsi souvent mieux les vedettes internationales que les héros et l'histoire de leur propre pays.

\souspartie{3. Une méconnaissance croissante du milieu local}
Parallèlement, la transmission traditionnelle recule : les veillées disparaissent au profit de la télévision et du téléphone, l'exode rural éloigne les jeunes de leur village, et de nombreux enfants ne parlent plus bien leur langue maternelle. Beaucoup ignorent l'histoire de leur village, les savoirs de leurs grands-parents, les noms des plantes médicinales ou les réalités économiques de leur région. Le constat de Mialaret est donc largement fondé.

\transition{Cette situation n'est pas sans conséquences pour les individus et pour la société.}

\partie{II. Deux conséquences possibles}

\souspartie{1. L'acculturation et la perte des repères}
Exposés en permanence aux modèles culturels venus d'ailleurs, de nombreux jeunes adoptent des modes de vie, des valeurs et des comportements étrangers, parfois au détriment des valeurs de respect, de solidarité et de dignité de leur société. La perte des langues nationales, des contes et des savoirs traditionnels appauvrit le patrimoine culturel et crée une crise d'identité : on risque de devenir étranger dans son propre pays.

\souspartie{2. Le déracinement et le désintérêt pour le développement local}
Celui qui ne connaît pas son milieu ne peut pas le transformer. Des jeunes formés selon des modèles extérieurs méprisent parfois les métiers de l'agriculture et de l'élevage, rêvent de partir en ville ou à l'étranger et se désintéressent des problèmes de leur localité. Les projets de développement conçus sans connaissance des réalités locales échouent souvent. Le village se vide de ses forces vives.

\transition{Pour éviter ces conséquences, il faut rechercher un équilibre entre l'ouverture au monde et l'enracinement dans le milieu local.}

\partie{III. Deux mesures pour trouver un équilibre}

\souspartie{1. Une école enracinée dans le milieu}
L'école doit enseigner, à côté des connaissances universelles, l'histoire, la géographie, les langues et les cultures locales. Les écoles bilingues, qui utilisent une langue nationale et le français, l'étude du milieu, les sorties pédagogiques, l'invitation d'anciens, d'artisans et de griots en classe permettent aux élèves de mieux connaître leur environnement. Les activités pratiques (jardin scolaire, élevage, artisanat) rapprochent l'école de la vie du village.

\souspartie{2. Mettre les médias et le numérique au service de la culture locale}
Les outils qui éloignent du village peuvent aussi en rapprocher : radios communautaires en langues nationales, émissions consacrées aux traditions, contenus numériques sur l'histoire et les cultures locales, festivals culturels (comme la Semaine nationale de la culture). Les familles ont aussi un rôle à jouer en parlant les langues nationales à leurs enfants et en perpétuant la transmission entre générations.

\partie{Conclusion}

La réflexion de Gaston Mialaret est fondée : alors qu'autrefois la connaissance se limitait au village, la société de l'information ouvre aujourd'hui sur le monde, mais au prix d'une méconnaissance croissante du milieu local. Cette situation peut entraîner l'acculturation des jeunes et leur déracinement, au détriment du développement local. Pour trouver un équilibre, l'école doit s'enraciner dans le milieu, et les médias comme les familles doivent valoriser les cultures locales. Il s'agit, selon la formule de Senghor, d'être à la fois « enraciné dans ses valeurs » et « ouvert aux apports » des autres : connaître son village pour mieux s'ouvrir au monde.

\end{spacing}
\end{document}
